Reference protection compares a proposed route with an executable alternative
before committing either route.  Let $P_b$ minimize the forecast matched
service score on the common nonempty portfolio $\mathcal P$.  Fix the same
deterministic ordering for every tie.  For each $P\in\mathcal P$, define
\begin{align}
\Delta(P,M)&=\widehat Q_t(P,M)-\widehat Q_t(P_b,M),\\
U(P)&=\max_{M\in\mathcal U_V}\Delta(P,M).
\label{eq:reference-difference}
\end{align}
The subtraction in \eqref{eq:reference-difference} uses the same matrix for
both routes.  It is not the difference between two separately maximized
absolute scores.  Let $P_u$ be the first minimizer of $U$ and set
\begin{align}
\mu&=\nu\max\{1,|\widehat Q_t(P_b,M_c)|\},\\
\tau&=10^{-10}\max\{1,\max_{P,M}|\widehat Q_t(P,M)|\}.
\end{align}
For a nonnegative margin ratio $\nu$, PC with reference protection returns
\begin{equation}
P_g=\begin{cases}
P_u,&U(P_u)<-\mu-\tau,\\
P_b,&\text{otherwise}.
\end{cases}
\label{eq:protected-route}
\end{equation}
Thus an insufficient predicted advantage retains the forecast matched route,
including its complete travel and completion schedule.  This rule changes
route selection only.  It does not change the common service allocations,
work orders, forecasts, uncertainty matrices, or execution constraints.
The comparison is made once at dispatch and uses no realized future outcome.
The margin is selected on validation observations and fixed before subsequent
evaluation.  The tolerance prevents near ties from replacing the reference
but is not a bound on all floating point errors in the forecast and score.

To separate the influence of uncertain parameters from route construction,
a second comparison uses six priorities that contain no transition
coefficient.  They are the training mean arrivals, the weighted current
threat, current threat multiplied by forecast exposure, total forecast
arrivals, total forecast energy plus backlog, and a uniform priority.
All selectors within this comparison receive the same resulting routes.
An observation supported score retains road persistence, spatial spread,
power persistence, and the fitted power to road association, while setting
communication persistence and the other four couplings to zero in the
score only.  Realized dynamics are unchanged.  This is an ablation of
unobserved terms, not an assertion that their physical effects vanish.
Changes between the two route portfolios cannot be attributed to score
selection alone.
